The corresponding functional block-diagram of the defended nonlinear system is
presented in \autoref{fig:defended-system}. The diagram extends the attacked
system (\autoref{fig:attacked-system}) by adding three redundant reference
channels and a voter block $V$. The reference signal $r(k)$ is broadcast to
three channels; one channel is corrupted by the adversarial perturbation
$\epsilon(k)$. Each channel computes its own tracking error $e_i(k)$. The
voter block receives all three error signals, computes the pairwise
disagreement $V_{ij}(k)$ from \autoref{eq:pairwise_disagreement}, selects
the minimum-variance pair via \autoref{eq:pair_selection}, and outputs the
defended error $\hat{e}(k)$. The remainder of the control path --- sign
activation $S$, learning rate $\beta$, nonlinear functions $f(x)$ and $g(x)$,
and NN solver $LAESS$ --- is identical to the nominal system in
\autoref{fig:clean-system}.

\begin{figure}[H]
    \centering
    \includesvg[width=0.9\textwidth]{defended-system}
    \caption{Functional block-diagram of the controlled nonlinear system with
    voting scheme defense.}
    \label{fig:defended-system}
\end{figure}
